\documentclass[ASNA,twocolumn]{USG}
% Wiley template downloaded from the Electronics Letters author guide (WileyDesign/).
% Build: latexmk (see latexmkrc; XeLaTeX, class and .bst taken from WileyDesign/).
% Spelling: British English (IET).
%
% Length (EL author guide, Letter): 2500-4200 words of main text; each figure or table
% counts as 800 words; references are not counted; captions concise.
% Plan: Fig. 1 + Fig. 2 + Table 1 = 2400 word-equivalents -> main text about 1500-1800 words.
% A Letter must: present original results; state novelty; compare with recently published
% approaches; show a clear engineering application.

\usepackage{booktabs}
\usepackage{xcolor}
\newcommand{\todo}{\textcolor{red}{TODO}}

% USG.cls draws the sample journal's logo (allergy.eps) on the title page; remove it.
\usepackage{etoolbox}
\makeatletter
\patchcmd{\@maketitle}{\includegraphics[width=56mm]{allergy.eps}}{}{}%
  {\PackageWarning{paper}{Could not remove the sample logo}}
\makeatother

% Draws a framed placeholder until the figure file exists.
\newcommand{\figorplaceholder}[1]{\IfFileExists{#1.pdf}{\includegraphics[width=\columnwidth]{#1}}%
  {\fbox{\parbox[c][40mm][c]{0.95\columnwidth}{\centering\detokenize{#1} (\todo)}}}}

\articletype{LETTER}
\journal{Electronics Letters}
\volume{0}
\copyyear{2026}
\startpage{1}
\articledoi{10.1002/0000}
\received{00}
\revised{00}
\accepted{00}
% Placeholders required by USG.cls; the publisher fills them in.

\makeatletter\@ifundefined{proposition}{\newtheorem{proposition}{Proposition}}{}\makeatother
\begin{document}

\title{Logic-gate complexity of digital chaotic systems}
\titlemark{Logic-gate complexity of digital chaos}% running title, < 40 characters (EL)
% Alternatives: "A gate-count law for fixed-point chaotic maps and flows"
%               "How many logic gates does chaos need?"

\author[1]{Erivelton Nepomuceno}[https://orcid.org/0000-0002-5841-2193]
\author[1]{Thalita Nazar\'e}[https://orcid.org/0000-0002-4372-6608]
\author[1]{Arthur M. Lima}[https://orcid.org/0000-0002-0542-3753]
\author[2]{Chi-Tsun Cheng}[https://orcid.org/0000-0003-3306-6148]
\authormark{NEPOMUCENO \textsc{et al.}}

\address[1]{\orgdiv{Hamilton Institute and Department of Electronic Engineering, }%
\orgname{Maynooth University, }\orgaddress{\state{Maynooth, }\country{Ireland}}}
\address[2]{\orgdiv{Department of Mechanical, Manufacturing and Mechatronic Engineering, School of Engineering, }%
\orgname{RMIT University, }\orgaddress{\state{Melbourne, }\country{Australia}}}
% Cheng's affiliation taken from his IEEE CAS contact page: confirm with him.

\corres{Erivelton Nepomuceno (\email{erivelton.nepomuceno@mu.ie})}

\fundingInfo{Taighde \'Eireann -- Research Ireland supported this research through grants
18/CRT/6049 and 21/FFP-P/10065, with additional support from Rinn AI (Research \&
Innovation in Data Science and AI) through grant 23/RC/13506. Arthur M. Lima
received an Electronic Engineering PhD co-funded scholarship, and the Conselho
Nacional de Desenvolvimento Cient\'ifico e Tecnol\'ogico (CNPq) supported
Erivelton Nepomuceno through grant 444154/2024-8.}

\keywords{chaos | digital circuits | field programmable gate arrays | fixed-point arithmetic | logic gates | random number generation}
% TODO: check whether EL requires keywords from a controlled list.

\abstract[ABSTRACT]{Chaotic generators for lightweight hardware have been chosen before
synthesis, yet their digital cost is reported case by case. This letter introduces the logic-gate
complexity $L(n)$ of chaotic maps and flows, predicted from their equations and design choices.
It serves to compare chaotic generators before synthesis.
Taking two-input logic gates as the primitive, we present the key design steps: the update
is written in $n$-bit fixed point as additions, piecewise-linear operations, products and
constant multiplications, each operation is replaced by an adder or multiplier circuit, and
the gates that drive an output are counted. We evaluate the technique on eight chaotic
systems and show that the 64-bit circuits preserve the main qualitative chaotic behaviour
of float64, with the largest Lyapunov exponent within 3.6\% and the state distribution
within 12\%, except for the Gingerbreadman map.}

\maketitle

%--------------------------------------------------------------------------
\section{Introduction}\label{sec:intro}
Hardware resource and energy constraints limit the implementation of chaotic
generators in lightweight and embedded devices.
Chaotic generators are research options for Internet-of-Things (IoT) communications
and image encryption~\cite{karagiorgos2024}, and lightweight stream
ciphers~\cite{ding2019}, motivating comparisons of candidate systems before
synthesis. FPGA studies usually report resources, speed and power for particular
systems, word widths and integration methods
\cite{mansingka2013, rezk2020, guillen2021, jauhari2023, damaj2024, babajans2025, alloun2025, tasdemir2025}, whereas others examine
finite-precision degradation~\cite{li2005, mendez2019}. Removing multipliers or changing
number representations reduces cost, but reported savings remain system-specific
\cite{rezk2020, nepomuceno2019, teixeira2019}.

Few equation terms or nonlinearities~\cite{sprott1994} do not necessarily imply low
circuit cost. Comparing five systems, Indrusiak et al.\ \cite{indrusiak2005} linked
the Linz--Sprott flow's lower FPGA cost to its multiplier-free nonlinearity.
Can gate cost be predicted from the equations and word width before synthesis?

We address this question with a gate-count law, illustrated on eight maps and flows,
that separates state products and non-dyadic constants from linear-cost operations,
using architecture-dependent coefficients. As an early-stage design indicator,
it identifies costs potentially removable by changing parameters or the system.
We compare gate counts with FPGA synthesis and estimate energy using a
switching-activity model.

%--------------------------------------------------------------------------
\section{Gate-level implementation}\label{sec:impl}
Each system has a single description that provides float64 reference results
or produces a gate-level netlist of two-input AND, OR and XOR gates and inverters. State variables use
$n$-bit two's-complement fixed-point numbers, with the number of integer bits set by
the float64 attractor range. Additions use ripple-carry adders, and
products use radix-4 Booth multipliers that keep only the columns needed for the exact
floor of the product. A constant is applied as a canonical-signed-digit chain of shifts
and additions with four guard bits, rounded once per linear combination, and is taken
from float64 except in Fig.~\ref{fig:Ln}. Flows are
integrated by forward Euler with a step $\mathrm{d}t = 2^{-k}$, so that multiplying by
$\mathrm{d}t$ requires only wiring, although the step can also be included in the constants.
Sine and
cosine take their argument in turns, fold it to a quarter turn and evaluate an odd
Taylor polynomial by Horner's rule, with the degree chosen for the word width.
During netlist construction,
operations on known constants are simplified and identical gates are shared,
counting only gates that contribute to an output.
Gate counts, fidelity metrics and energy estimates for 73 maps and flows
(117 implementations, including RK4 versions of the flows) are available in the
repository given in the Data Availability Statement; this letter reports
representative systems. Across 200 states from each
attractor, the netlist and the integer model used for the long runs
gave the same output bits in every check.

%--------------------------------------------------------------------------
\section{Gate-count law}\label{sec:law}
State products and constant multiplications account for the quadratic cost.
A product of two $n$-bit signals uses a radix-4 Booth multiplier whose gate count
grows as $23n^2/4$, regardless of the number of integer bits, giving $a = 5.75$.

Multiplication by a constant $c$ uses a chain of ripple-carry adders, one per non-zero
canonical-signed-digit (CSD) digit, at five gates per bit per adder. When $c$ is
stored at the word resolution, it has $\rho_c n$ non-zero digits, where $\rho_c$ is its
digit density. CSD has no adjacent non-zero digits, so $\rho_c \le 1/2$. Dyadic
constants have a fixed number of digits and cost $O(n)$. Multiplication by
$\mathrm{d}t = 2^{-k}$ requires only wiring, and added constants cost $O(n)$. The gate count is not a property of the equations alone: it also depends on design
choices, namely how the update is written, the adder and multiplier algorithms, the
discretisation scheme and step, and the word width. The construction is therefore stated
as a procedure. (i) Write the map, or the Euler update of the flow with
$\mathrm{d}t=2^{-k}$, as a sequence of additions, piecewise-linear operations, products
and constant multiplications. (ii) Choose $n$ and the integer bits from the attractor
range. (iii) Replace each addition by a ripple-carry adder, each product by a Booth
multiplier, each constant by its CSD shift-and-add chain and each dyadic factor by
wiring. (iv) Simplify the netlist and count the gates that drive an output.
For this procedure:

\begin{proposition}\label{prop:law}
Assume ripple-carry adders, radix-4 Booth multipliers, a fixed number of integer bits and
constants stored to the word resolution. Write the update as a straight-line program of
additions, piecewise-linear operations, $M$ products of two state-dependent signals and
$C$ multiplications of a signal by a non-dyadic constant $c$ (not of the form $k/2^m$).
Then the gate count per iteration satisfies
\begin{equation}
  L(n) = 5.75\,M\,n^2 + 5\sum_{c=1}^{C} \rho_c\, n^2 + O(n)
       \le \left(aM + bC\right) n^2 + O(n),
  \label{eq:law}
\end{equation}
with $a = 5.75$ and $b = 2.5$, where $\rho_c \le 1/2$ is the CSD digit density of $c$.
\end{proposition}
\begin{proof}[Sketch]
Each Booth product costs $23n^2/4+O(n)$ gates, each constant costs one $5n$-gate adder per
non-zero CSD digit, of which there are $\rho_c n+O(1)$ with $\rho_c\le1/2$, and every other
operation costs $O(n)$.
\end{proof}
Because $M$ and $C$ depend on how the update is written ($xy+xz$ or $x(y+z)$), we take the
form with the fewest products; Eq.~(\ref{eq:law}) is then an upper bound on the minimal
gate count, and its growth order, not the exact count, defines the cost class. Fits to the measured counts give
the predicted quadratic terms, including $5.75n^2$ for the logistic map ($M = 1$, $C = 0$).
The R\"ossler flow, with $z(x-c)$ computed as one product, has one state product and
one constant $a\,\mathrm{d}t = 0.2\,\mathrm{d}t$ with $\rho = 1/2$, giving $8.25n^2$.
The Lorenz flow, with $x(28-z)$ as one product and dyadic $\sigma\,\mathrm{d}t = 10/2^7$, has
two state products and one constant $\beta\,\mathrm{d}t = 1/48$ with $\rho = 1/2$, giving $14.0n^2$
(Table~\ref{tab:systems}).

Sine, cosine and exponential functions use polynomials whose degree grows with $n$,
adding a cubic term. This gives three classes by operation: piecewise-linear systems with $M = C = 0$ and
linear cost (Gingerbreadman, Lozi, tent), polynomial systems with state products or
non-dyadic constants and quadratic cost (MDCS, logistic, Chua, R\"ossler, Lorenz), and
transcendental systems with cubic cost, such as the sine map
($\mu=0.99$, about $0.27n^3+82n^2$, 293\,018 gates at 64 bits; see the repository). Figure~\ref{fig:Ln} shows the systems of
Table~\ref{tab:systems}. The polynomial class includes Chua's circuit, whose
piecewise-linear nonlinearity uses non-dyadic constants.
% Class counts from logical-chaos.md (71 systems with IDs) plus the two degenerate maps.
% Constants in the code come from float64 (53 significant bits): beyond n ~ 54 their CSD
% digits saturate at about 27 and the constant term becomes linear (~135n per constant);
% this is why Chua's 16-64 fit gives only 1.44 n^2. Fig. 1 uses exact constants.

\begin{figure}[t]
\centering
\figorplaceholder{fig1_Ln}
\caption{Gates per iteration against word width for systems in the piecewise-linear
and polynomial classes, with constants stored to the word resolution. Dashed lines
show slopes 1 and 2.}
\label{fig:Ln}
\end{figure}
% Fig. 1 from analysis/fig1_data.py + analysis/figures.py (ripple + Booth, exact constants).

%--------------------------------------------------------------------------
\section{Design rule}\label{sec:rule}
Avoiding state products and choosing dyadic constants gives low-cost circuits.
Piecewise-linear operations such as $|x|$, $\min$ and saturation keep linear cost,
so Gingerbreadman, Lozi and tent maps need fewer than 1000 gates at 64 bits.
For the tent map, float64 $\mu=1.99$ saturates at 16 CSD digits from about 56 bits,
costing $84n+195$ gates. Choosing $\mu=2-2^{-7}$ needs two digits and
$14n-17$ gates, reducing the 64-bit count from 5571 to 879 by a factor of 6.3
while the Lyapunov exponent $\lambda=\ln\mu$ changes from 0.6881 to 0.6892.
For the logistic map, $r=63/16$ removes quadratic constant-multiplication cost,
but the state product retains $5.75n^2$.

We call the potential reduction towards a comparable linear-cost design the
headroom. Equation~(\ref{eq:law}) identifies its product and non-dyadic-constant
contributions, with the tent example removing 84\% of the 64-bit gates. This is an
architecture-relative screening indicator, not an absolute silicon-cost prediction
or a universally attainable floor. Removing an essential product changes the system,
and dyadic parameter substitutions must preserve the required dynamics, while
transcendental functions retain their cubic cost unless replaced.

Replacing $x^2$ with $|x|$ also motivates comparing logistic with tent and H\'enon
with Lozi. At 64 bits, tent needs 879 gates versus 23\,705 for logistic, and
Lozi ($a=7/4$, $b=1/2$) needs 754 versus 40\,846 for H\'enon ($M=1$, $C=2$).
Their respective Lyapunov exponents are 0.69 versus 0.53 and 0.51 versus 0.42,
so lower cost does not imply a lower exponent in these examples.

%--------------------------------------------------------------------------
\section{Validation and comparison}\label{sec:validation}
Eight 64-bit netlists were synthesised with Vivado 2025.2 for a Zynq UltraScale+
xczu7ev, using out-of-context synthesis, place and route, and one gate per lookup
table (LUT). Mapping one gate per LUT (hierarchy kept, no LUT combining) isolates the
netlist cost, so this comparison checks the netlists rather than predicting
LUT counts; default synthesis, which packs gates into six-input LUTs and maps
products to carry chains and DSP blocks, gives % TODO: HDL results
\todo. LUT counts differ from gate counts by $-8\%$ to $+14\%$, with
a log--log slope of 0.99 ($R^2=0.999$) over two orders of magnitude
(Table~\ref{tab:systems}), preserving the separation between classes.

Four-bit carry-lookahead adders give $(a,b)=(9.125,4.75)$, array multipliers give
$(8.5,2.5)$, and combining both gives $(15.25,4.75)$. Equation~(\ref{eq:law}) predicts
the Lorenz, R\"ossler and logistic quadratic coefficients in each architecture.
Slopes for the piecewise-linear class remain 1.00--1.02 over 64--128 bits, with bitwise agreement
between every netlist and its reference.

The switching-activity model assigns 40\,fJ per LUT toggle and 10\,fJ per flip-flop
per cycle. For 114 of the 117 64-bit circuits in the repository,
estimated energy per output bit scales with gates per state bit ($L(64)$ over the number
of state bits), with slope 0.91 and
$R^2=0.94$ (Fig.~\ref{fig:energy}). Estimates are 0.06\,pJ for Gingerbreadman
and 2.4\,pJ for Lorenz. The piecewise-linear class has the lowest estimated
energy per output bit under this model.

Published FPGA comparisons require care because devices, widths and mappings differ.
Chua used 5544 LUTs versus 2268 for Lorenz in~\cite{indrusiak2005}, despite its
piecewise-linear nonlinearity, and Eq.~(\ref{eq:law}) identifies costs from its three
non-dyadic constants. Multiplierless Chua and modified Lorenz and R\"ossler generators with
power-of-two constants need 178--224 LUTs at 20 fractional bits~\cite{rezk2020},
consistent with linear cost.

LUT counts omit multiplication mapped to digital signal processing (DSP) blocks.
The 32-bit Lorenz and R\"ossler circuits in~\cite{jauhari2023} use 378 and 500 LUTs
plus 8 and 4 DSP blocks, respectively. Reported counts of 2748 LUTs for an 8-bit
Euler core and 5 for RK4~\cite{damaj2024} likewise resist comparison by integration
order alone. Gate counts provide a common basis before synthesis.

\begin{table}[t]
\caption{Representative systems showing products between state variables $M$, non-dyadic constants
$C$, fitted gate count per iteration, gates at 64 bits and FPGA LUTs.\label{tab:systems}}
\centering
\footnotesize
\setlength{\tabcolsep}{3pt}
\begin{tabular}{@{}lcclrr@{}}
\toprule
System & $M$ & $C$ & $L(n)$ & $L(64)$ & LUTs\\
\midrule
Gingerbreadman      & 0 & 0 & $7n + 2$                & 450     & 450\\
Lozi~\cite{lozi1978} & 0 & 0 & $12n - 14$              & 754     & 758\\
Tent$^{\rm a}$      & 0 & 0 & $14n - 17$              & 879     & 915\\
MDCS$^{\rm b}$      & 1 & 0 & $1.50n^2 + 4.7n - 26$   & 6432    & 6838\\
Logistic            & 1 & 0 & $5.75n^2 + 2.5n - 7$    & 23\,705 & 23\,994\\
Chua$^{\rm c}$      & 0 & 3 & ---$^{\rm c}$ & 22\,647 & 25\,724\\
R\"ossler~\cite{rossler1976} & 1 & 1 & $8.22n^2 - 24n - 237$   & 31\,925 & 29\,704\\
Lorenz~\cite{lorenz1963} & 2 & 1 & $14.0n^2 - 51n - 281$   & 53\,799 & 49\,450\\
\bottomrule
\end{tabular}
\par\smallskip
\parbox{\columnwidth}{\scriptsize LUT counts use 64-bit netlists, one gate per LUT, Vivado 2025.2,
Zynq UltraScale+ xczu7ev (ZCU104).
$^{\rm a}$ $\mu = 2 - 2^{-7}$.
$^{\rm b}$ Minimal digital chaotic system~\cite{nepomuceno2019} with a dedicated digit-complement multiplier.
$^{\rm c}$ Float64 constants, whose CSD digits saturate within 16--64 bits, so no stable
quadratic fit exists; the cost is close to linear beyond 64 bits.}
% Parameters removed from the notes (all in the text or in logical-chaos.md): Lozi a = 7/4,
% b = 1/2; logistic r = 63/16; Chua alpha = 15.6, m0 = -8/7, m1 = -5/7; Rossler a = b = 0.2,
% c = 5.7 (c and b are free); Lorenz sigma = 10, rho = 28, beta = 8/3.
\end{table}
% M and C read from the code in ec_artifact_src/ec_artifact/compute/zoo.py: dt is folded into
% the constants; only constants that multiply a signal count (additive ones are free);
% integers such as 10, 28 have short CSD chains (linear cost).
% Measured with BitB (analysis/const_cost.py, analysis/const_exact.py), n = 16..256:
%   one product: exactly 5.75 n^2 for any number of integer bits I (a = 23/4);
%   one exact constant: 5*rho*n^2, rho = CSD digit density <= 1/2 -> b <= 2.5
%     (1/5, 8/3: rho = 1/2, coef 2.500; 15.6*(2/7)*2^-7: coef 1.69);
%   float64 constant: CSD digits saturate at up to about 27, so cost up to ~ 5*27*n = 135n
%     beyond n ~ 54; tent mu = 1.99 saturates at 16 digits from n = 56: exactly 84n + 195.
% Law check, L ~ (5.75 M + 2.5 C) n^2 (16-64 fits):
%   Logistic 5.75 (fit 5.75) OK; Lorenz 14.0 (fit 14.0) OK; Rossler 8.25 (fit 8.25) OK;
%   Chua 7.5 vs fit 1.44: three float64 constants with rho < 1/2 that saturate within
%     16-64 bits (I = 4, so F is large) -> constant term is not n^2 over the fit window;
%   MDCS 5.75 vs fit 1.50: special multiplier, outside the architecture.
% LUTs from the authors' Vivado runs (2026-11-05).
% pJ/bit column removed (2026-11-05); energy only in Sec. 5 and Fig. 2.

\begin{figure}[t]
\centering
\figorplaceholder{fig2_energy}
\caption{Estimated energy per output bit against gates per state bit for 114 of the
117 64-bit circuits in the repository. Markers show the class, the line shows the
least-squares fit, and labelled points refer to the systems in Table~\ref{tab:systems}.}
\label{fig:energy}
\end{figure}

%--------------------------------------------------------------------------
\section{Fidelity}\label{sec:fidelity}
A low-cost circuit is useful only if it preserves the chaotic behaviour of interest.
For each system of Table~\ref{tab:systems}, we compared the 64-bit circuit with float64
over $10^5$ iterations for maps and 500 time units for flows, using the same
integration method and step, and compared the largest Lyapunov exponent and the
distribution of $x$ (Wasserstein distance) with the difference between two float64 runs
started from nearby initial conditions.
The relative error of the exponent is $1.1\times10^{-8}$ for the tent map with
$\mu = 2 - 2^{-7}$, $6\times10^{-5}$ for MDCS, 0.2\% for Lozi, 0.9\% for the logistic map,
1.6\% for Chua and Lorenz, and 3.6\% for R\"ossler, whose exponent and distribution both
lie within the float64 noise floor. The distribution error stays below 0.5\% for
Lozi, tent, MDCS, logistic and R\"ossler, but reaches 9\% for Lorenz and 12\% for Chua.
The Gingerbreadman map is the exception, with a 28\% exponent error and a 3\% distribution
error, so its low cost comes with reduced fidelity at 64 bits.
Long-term periodicity of the digital maps is outside the scope of this study \cite{li2005}.

%--------------------------------------------------------------------------
\section{Conclusion}\label{sec:conclusion}
For this architecture, each state product contributes $5.75n^2$ gates and each
non-dyadic constant at most $2.5n^2$, whereas dyadic constants and piecewise-linear operations
retain linear cost. Other adders and multipliers change the coefficients but preserve
this form, and LUT counts for eight synthesised systems agree within 14\% of gate counts.
Piecewise-linear maps need a few hundred gates
and estimated 0.06--0.14\,pJ per output bit.
The law supports screening before synthesis by identifying architecture-relative
headroom, whose realisation depends on preserving the required dynamics.

\bmsubsection*{Author Contributions}
\textbf{Erivelton Nepomuceno:} Conceptualization, Methodology, Software, Investigation,
Writing -- original draft, Resources, Funding acquisition.
\textbf{Thalita Nazar\'e:} Methodology, Software, Investigation, Writing -- original draft.
\textbf{Arthur M. Lima:} Methodology, Software, Investigation, Writing -- original draft.
\textbf{Chi-Tsun Cheng:} Conceptualization, Writing -- review \& editing.

%\bmsubsection*{Acknowledgments}
%\todo.

\bmsubsection*{Conflicts of Interest}
The authors declare no conflicts of interest.

\bmsubsection*{Data Availability Statement}
The code and data supporting this study, including gate counts, fidelity metrics and
energy estimates for all 73 systems, are openly available in Zenodo at
https://doi.org/\todo.
% TODO: check which statements EL requires (data availability, funding, ORCID).

\bibliography{refs}

\end{document}
